\documentclass[11pt]{article}
\usepackage[margin=1.05in]{geometry}
\usepackage{amsmath,amssymb}
\usepackage{booktabs}
\usepackage{enumitem}
\usepackage[hidelinks]{hyperref}
\usepackage{xcolor}

\newcommand{\Jtraj}{\mathbf{J}_{\mathrm{TRAJ}}}
\newcommand{\Jdemo}{\mathbf{J}_{\mathrm{DEMO}}}
\newcommand{\Ciss}{C_{\mathrm{ISS}}}
\newcommand{\kstar}{k^{\star}}

\title{Segment-Dependent Action Chunk Length via a Video Stability Predictor\\
\large Complete Implementation Document, v0.2}
\author{Aryan Goyal}
\date{July 2026 (status and results section added July 20)}

\begin{document}
\maketitle

\begin{abstract}
\noindent
This document describes the full plan for a system that changes the executed action chunk
length of an imitation-learned robot policy at run time. The decision signal is a small
video model that watches the last 16 camera frames and predicts whether the robot is
currently in an open-loop stable segment or an open-loop unstable segment. In stable
segments the policy executes a long chunk of actions without feedback. In unstable
segments the policy replans at every step. The ground-truth stability labels come from a
simulator, using a standard perturb-and-replay procedure that measures the finite-time
Lyapunov exponent. Only one small network head is trained. The policy and the video
encoder are reused without modification. This document explains the models, the labeling
procedure, the datasets, the training recipes, the runtime controller, the experiments,
the baselines, the risks, and the compute budget.
\end{abstract}

\tableofcontents
\newpage

%=====================================================================
\section{The problem in plain terms}

An imitation-learned policy such as Diffusion Policy or ACT predicts a chunk of future
actions at once. The system designer must choose how many of those predicted actions to
actually execute before asking the policy again. Call this number $k$, the executed chunk
length. If $k=1$ the policy replans at every step, which is fully closed loop. If $k$ is
large the robot runs open loop for $k$ steps between replans.

The common belief in the literature is that small $k$ is always safe and large $k$ is
only a matter of efficiency. Recent theory shows this belief is wrong. The correct choice
of $k$ depends on a property of the environment, not of the policy:

\begin{itemize}[leftmargin=2em]
  \item In \textbf{open-loop stable} segments (free-space reaching and transport, where
        the low-level controller pulls small errors back to the reference), the chunk
        must be \emph{longer} than a threshold $\kstar$. Short chunks are actively
        harmful here, because every replan injects the policy's estimation error into
        the system, and per-step replanning injects it at every step.
  \item In \textbf{open-loop unstable} segments (contact events such as grasping and
        insertion, where small errors grow), long open-loop execution lets errors
        compound without correction. The chunk must be \emph{short}, and per-step
        feedback is the right mode.
\end{itemize}

A long-horizon task contains both kinds of segments. Therefore no single constant $k$
can be correct for the whole task. The hypothesis of this project is that $k$ should be
switched per segment, and that the switching signal should measure the stability of the
environment rather than the confidence of the policy.

\subsection{The theory behind the two constraints}

The relevant results are from Simchowitz et al.\ (arXiv:2503.09722) and Zhang et al.\
(arXiv:2507.09061). Let $\Jtraj$ be the trajectory error of the learned policy at
execution time and $\Jdemo$ its regression error on the demonstration distribution.
Exponential error compounding means
\begin{equation}
\Jtraj \;\gtrsim\; C^{T}\,\Jdemo, \qquad C>1,
\end{equation}
where $T$ is the horizon. The stability notion is exponential incremental input-to-state
stability (EISS): a system is $(\Ciss,\rho)$-EISS with $\rho\in(0,1)$ if perturbations
between two trajectories contract geometrically over time.

The theory splits into two cases:
\begin{enumerate}[leftmargin=2em]
  \item \textbf{Stable dynamics, closed-loop execution still fails.} Even when the
        dynamics contract, a smooth Markovian policy that replans every step suffers
        exponential compounding. The feedback channel is the mechanism: the policy
        carries estimation error, and closing the loop injects that error at every step.
  \item \textbf{Chunking fixes the stable case, if the chunk is long enough.} Executing
        $k$ actions open loop means error is injected once per chunk instead of once per
        step, and the contraction absorbs it between injections. This works only above a
        threshold:
        \begin{equation}
        k \;>\; \kstar \;=\; \log(1/\rho)^{-1}\cdot\log\bigl(\mathrm{poly}(L_\pi,\Ciss)\bigr),
        \end{equation}
        where $L_\pi$ is the policy Lipschitz constant. Above $\kstar$ the trajectory
        error is horizon-free, $\Jtraj \leq O_\star(1)\,\Jdemo$. Below $\kstar$ the
        exponential lower bound still applies.
  \item \textbf{Unstable dynamics cannot be fixed at execution time.} For open-loop
        unstable dynamics, exponential compounding is unavoidable for any policy trained
        only on expert data. The admissible fix is on the data side: injecting noise
        into the demonstrations (Theorem 2 of Zhang et al.) reduces the error growth
        from exponential to polynomial in $T$. In these segments short chunks and
        per-step feedback are the best execution-side option, and chunking is known to
        degrade performance empirically.
\end{enumerate}

Applied segment by segment, the two constraints have opposite signs: contracting segment
$i$ requires $k \geq \kstar_i$, while expansive segment $j$ requires $k \lesssim
1/\lambda_j$ where $\lambda_j$ is the local expansion rate. A constant $k$ generically
violates one of them.

\subsection{Why existing adaptive-chunking methods do not solve this}

Existing adaptive methods (AAC, DVAC, HiPolicy, BID, AQC, competency-adaptive execution)
all condition on policy-side signals: action entropy, denoising variance, resampling
coherence, or learned value estimates. These signals can only ever license commitment.
They say that committing is permissible when the policy is confident. They cannot say
that committing is \emph{required}. The harm of not committing in a stable segment lives
in the feedback channel, which is a property of the closed-loop transition map and is
absent from the policy's output distribution. A policy can be confident yet wrong in a
contact segment, and uncertain in a benign free-space segment. Our signal measures the
environment instead, so it can produce the minimum-length constraint that policy-side
signals cannot produce even in principle. The disagreement cases between the two kinds
of signal are the natural experiment, and we test them explicitly (Section
\ref{sec:experiments}).

%=====================================================================
\section{The models and how they fit together}
\label{sec:models}

The system has three model components with strictly separated roles. Only one of them is
trained in this project.

\subsection{The policy model (reused, never retrained)}

The policy is a standard chunking imitation policy:

\begin{itemize}[leftmargin=2em]
  \item \textbf{Stage 1:} Diffusion Policy, from the public real-stanford repository,
        with the released robomimic image configurations. We use released checkpoints
        where available, or retrain from the released config, which takes hours on one
        GPU.
  \item \textbf{Stage 2:} ACT on the ALOHA simulated tasks, to show the result is not
        specific to one policy class.
\end{itemize}

The key design decision is that the executed chunk length $k$ is an inference-time knob.
The policy always predicts its full action chunk; we only vary how many predicted actions
are executed before replanning. Because of this, the policy is never retrained, its loss
is never modified, and every comparison in the project runs the exact same policy
weights. Any performance difference between chunking strategies is therefore attributable
to the execution rule alone.

\subsection{The video model (frozen encoder plus a small trained head)}

The trained component is a stability predictor. Its only job is: given the recent video,
say whether the current segment is open-loop stable or unstable.

\begin{itemize}[leftmargin=2em]
  \item \textbf{Input:} the last 16 camera frames, plus proprioception (joint readings).
  \item \textbf{Backbone:} V-JEPA 2 ViT-L, completely frozen (public checkpoints,
        \texttt{facebook/vjepa2-*}). The reason for a video backbone rather than an
        image backbone is that stability is a dynamics property. Contact, incipient
        slip, and constrained motion are motion cues. They are visible in a clip and not
        in a single frame. V-JEPA 2 is pretrained by masked latent video prediction,
        which is exactly a motion-predictive objective, so its features are suited to
        this readout.
  \item \textbf{Head:} attentive pooling followed by a 2-layer MLP, which is V-JEPA's
        standard attentive-probe protocol, with reused code. Proprioception passes
        through a small MLP and is concatenated after pooling. This head is the only
        trained component in the entire project.
  \item \textbf{Outputs:} the probability that the current segment is open-loop
        unstable, $p(\text{OL-unstable})$; an auxiliary regression of the open-loop
        exponent $\hat\lambda_{\mathrm{ol}}$; and optionally the closed-loop pair
        $p(\text{CL-unstable})$ and $\hat\lambda_{\mathrm{cl}}$.
  \item \textbf{Ablations:} a single-frame DINOv2 backbone with the same head, to prove
        that the temporal information matters; and LoRA on the last encoder blocks in
        case the frozen probe underfits.
\end{itemize}

\subsection{The world model (label source, not a runtime component)}

The word ``world model'' plays two roles at two different times, and it is important to
keep them separate:

\begin{itemize}[leftmargin=2em]
  \item \textbf{Now, in simulation:} the ``world model'' is the MuJoCo simulator itself,
        used as a ground-truth oracle. It is not learned and it involves no gradients.
        We use \texttt{get\_state} and \texttt{set\_state} to reset the simulator to any
        recorded moment, perturb the first action, replay forward, and measure how fast
        trajectories diverge. This procedure generates the stability labels
        (Section \ref{sec:labels}).
  \item \textbf{Later, on a real robot:} there is no simulator to reset and replay, so a
        learned action-conditioned world model (V-JEPA 2-AC or DINO-WM) replaces the
        simulator as the label source. The world model imagines paired rollouts in its
        latent space: same actions, perturbed first action, and the label is the slope
        of the log latent divergence between the imagined rollouts. The labeling logic,
        the predictor architecture, and the losses are unchanged. Only the oracle that
        produces the labels is swapped. Hardware perturb-and-replay on a subsample of
        states is the backup option.
\end{itemize}

The world model is never in the control loop. At run time the only extra computation is
the frozen encoder and the small head.

\subsection{How the signal chooses the chunk length}
\label{sec:controller}

At each replanning point the controller runs the encoder and head on the last 16 frames,
reads out the predicted regime, and sets the executed chunk length:

\begin{center}
\begin{tabular}{@{}lll@{}}
\toprule
Predicted regime & Executed chunk length & Meaning \\
\midrule
OL-stable & $k = k_{\mathrm{long}}$ & run open loop; standard chunk \\
 & & (Diffusion Policy $T_a = 8$ to $16$, ACT $100$) \\
OL-unstable and CL-stable & $k = 1$ & per-step feedback; Markovian mode \\
Neither (both unstable) & $k = 1$ plus a flag & no execution-side fix exists here; \\
 & & the remedy is data-side noise injection \\
\bottomrule
\end{tabular}
\end{center}

Details of the rule:
\begin{itemize}[leftmargin=2em]
  \item \textbf{Hysteresis:} the controller requires two consecutive agreeing
        predictions before it switches regime, to avoid chattering at segment
        boundaries.
  \item \textbf{MVP:} the minimum viable version collapses to the first two rows, a
        binary switch on the open-loop head, since the third case is rare in these
        benchmarks. Both heads are trained from the start, so the full three-case rule
        costs nothing extra.
  \item \textbf{No derived $\kstar$:} we do not compute the theoretical threshold. The
        empirical $\kstar$ per segment comes from the Stage 1 sweep, defined as the
        smallest $k$ at which success plateaus in stable segments.
  \item \textbf{Continuous rule as ablation:} the auxiliary $\hat\lambda$ regression
        gives a continuous rule $k \approx c/|\hat\lambda|$ for free. This is an
        ablation, not the main method.
\end{itemize}

\subsection{The full pipeline in one line}

Simulator oracle $\rightarrow$ per-timestep stability labels $\rightarrow$ supervised
training of (frozen video encoder + small head) $\rightarrow$ at run time the predictor
watches recent frames $\rightarrow$ binary switch of the executed chunk length.

%=====================================================================
\section{Labeling open-loop and closed-loop stability}
\label{sec:labels}

This section describes exactly how the ground-truth labels are produced. There is no
learning anywhere in this procedure. It is a script, not a model.

\subsection{The quantity being measured}

The ground-truth quantity is the finite-time Lyapunov exponent (FTLE), the textbook
estimator of local exponential divergence (Benettin et al.\ 1980; used for reinforcement
learning policy robustness in arXiv:2410.10674). The question the label answers is: if a
small error of realistic size enters the system right now, does it grow over the next
segment or does it get absorbed?

The perturbation size is chosen to make this question exact. We set $\sigma_u$ equal to
the trained policy's validation action RMSE, so the probe injects exactly the size of
error the policy will actually make. With this choice the label directly measures the
error-compounding gain from the theory: whether an injected error of realistic size is
amplified or absorbed over the segment horizon.

The theory papers themselves (arXiv:2503.09722, arXiv:2507.09061) provide no empirical
classification procedure. They assume EISS and build synthetic systems that are stable
by construction. The per-segment labeling protocol below is ours, but every ingredient
is standard and citable: the FTLE, contraction analysis (Lohmiller and Slotine 1998),
and the imitation-learning community's open-loop replay test.

\subsection{The open-loop label (the primary label)}

The open-loop label is a property of the plant: the robot, its low-level controller, and
the environment together. For each demo timestep $t$, subsampled every 2 to 5 steps, we
run the following in MuJoCo:

\begin{enumerate}[leftmargin=2em]
  \item Reset the simulator to the recorded state at time $t$.
  \item \textbf{Nominal branch:} replay the recorded actions $a_{t:t+K}$ exactly. The
        final recipe uses $K = 24$; how $K$ is chosen is explained in Section
        \ref{sec:choosingK}.
  \item \textbf{Perturbed branches:} run $N = 8$ additional branches, identical to the
        nominal branch except that Gaussian noise with standard deviation $\sigma_u$ is
        added to the \emph{first action only}. All later actions are replayed unchanged.
        The robot does not react in any branch; the actions are fixed.
  \item Measure $d(\tau)$, the state-space distance between the perturbed branches and
        the nominal branch at each step $\tau$, using end-effector pose plus object
        pose. Fit
        \begin{equation}
        \lambda_{\mathrm{ol}}(t) = \text{OLS slope of } \log d(\tau) \text{ versus } \tau.
        \end{equation}
  \item Threshold with a deadband:
        \begin{equation}
        y_{\mathrm{ol}}(t) =
        \begin{cases}
        \text{unstable} & \lambda_{\mathrm{ol}}(t) > +\delta \\
        \text{stable} & \lambda_{\mathrm{ol}}(t) < -\delta \\
        \text{inherit from neighbor} & \text{otherwise.}
        \end{cases}
        \end{equation}
        Then apply a median filter over about 3 consecutive label stamps. The result is
        contiguous stable and unstable segments. Simulator contact flags are logged as
        an independent sanity channel, never as a label.
\end{enumerate}

The intuition is simple. If the nudged replays drift away from the clean replay, errors
grow on their own and the segment is unstable. If the nudged replays converge back to
the clean replay, the plant absorbs errors and the segment is stable.

\subsection{The closed-loop label (secondary, diagnostic)}

The closed-loop label uses the same perturbed pairs, with one change: in the perturbed
branch the trained policy acts from its own observations and replans at every step,
instead of replaying fixed actions. This yields $\lambda_{\mathrm{cl}}(t)$ and
$y_{\mathrm{cl}}(t)$ by the same fitting and thresholding procedure.

The comparison of $\lambda_{\mathrm{cl}}$ against $\lambda_{\mathrm{ol}}$ per segment
verifies the mechanism from the theory: in contracting segments, feedback injects
estimation error and should make things worse relative to open-loop replay; in expansive
segments the ordering should flip. The closed-loop head also enables the third row of
the runtime rule.

\subsection{How the probe window length K is chosen}
\label{sec:choosingK}

$K$ is the number of steps we watch the error for after the nudge. It is not
arbitrary: it is bracketed from below by two requirements and from above by two
others, and then fixed empirically inside that bracket.

\paragraph{Lower bound 1: the fit must see past the controller transient.}
$\lambda$ is the OLS slope of $\log d(\tau)$ against $\tau$. The first 0.2 to 0.5
seconds after the nudge (4 to 10 steps at 20 Hz) are dominated by the low-level
tracking controller settling, regardless of the regime. If $K$ barely exceeds this
transient, the fitted slope measures the transient rather than the plant regime.
$K$ must leave enough points beyond it for a stable line fit.

\paragraph{Lower bound 2: the window must cover the decision it certifies.}
The label's runtime meaning is ``it is safe to commit a chunk of actions open loop
starting here.'' A verdict measured over a window shorter than the chunk would not
certify the commitment being made. So $K$ must be at least the longest chunk the
controller may execute: $K \geq k_{\mathrm{long}} = T_p = 16$ for Diffusion Policy
(0.8 seconds at 20 Hz).

\paragraph{Upper bound 1: locality.}
The label is a per-segment quantity. If $K$ approaches the typical dwell time of a
segment, the window straddles regime boundaries and the fitted $\lambda$ blends a
contracting stretch with an expansive one. The median filter cannot repair a window
that is itself mixed. $K$ must therefore stay well below segment length.

\paragraph{Upper bound 2: the fit must stay in the exponential regime.}
The estimator assumes $d(\tau) \approx \sigma_u e^{\lambda \tau}$, which holds only
while the perturbation stays small relative to the workspace and no discrete event
(a new contact, a constraint activation) redirects the divergence. Over long windows
$d(\tau)$ saturates and $\log d$ flattens, biasing $\lambda$ toward zero exactly at
the unstable stamps we most need to detect. A larger $\sigma_u$ reaches saturation
sooner, so this bound couples $K$ to $\sigma_u$: with $\sigma_u$ set to the policy
RMSE, $K$ cannot be grown freely.

\paragraph{The empirical pick inside the bracket.}
The admissible range for our setup is roughly $16 \leq K \lesssim 30$. A six-variant
probe study on 30 Square demos selected $K = 24$ (1.2 seconds) with position-only
perturbation: 36.6\% of fixture-contact stamps sit above the deadband versus 12.7\%
at $K = 16$, while free-space false positives stay at 1.2\%. The Stage 0 synthetic
system provides the correctness check: at these settings the labeler recovers the
true $\lambda$ of a known piecewise-linear system to four decimal places.

$K$ is a per-embodiment constant, set by the controller settling time and the chunk
horizon, not a per-task tune. It is re-derived once for ALOHA in stage 2 (different
controller, different control rate), and the label hyperparameters $N$, $K$,
$\sigma_u$, $\delta$ get a sensitivity ablation in stage 2.

\subsection{Which states get labeled (coverage)}

Labels are generated along two kinds of trajectories:

\begin{enumerate}[leftmargin=2em]
  \item \textbf{Demonstration states:} every demo trajectory, subsampled as above. The
        replayed actions are the demo's recorded actions.
  \item \textbf{Policy rollout states, including failed rollouts:} trajectories produced
        by running the trained policy. Here the replayed actions in the nominal branch
        are the policy's own executed actions.
\end{enumerate}

The reason for including rollouts is distribution match. The predictor will run on the
states the policy actually visits, which include states outside the clean demo tube,
especially during recoveries and failures. Labeling only demos would train the predictor
on a narrower distribution than it faces at run time.

\subsection{Segmentation is an output, not an input}

We never segment trajectories by hand or by semantic phase. All datasets are used as
continuous, unannotated trajectories. The per-timestep FTLE labels, after thresholding
and median filtering, \emph{are} the segments. Stability regime is the only segment
notion the hypothesis uses. Whether the stability segments happen to align with semantic
phases such as reach, grasp, and insert is a qualitative result we report, not an
assumption we rely on.

\subsection{Cost}

The labeler runs pure simulator rollouts with no gradients, and every labeled timestep
is independent, so the work is embarrassingly parallel. The total is about 200 demos
plus a similar number of policy rollouts, times one label every 2 to 5 steps, times
$(N+1) = 9$ branches of $K = 16$ steps each. This is hours on CPU workers.

%=====================================================================
\section{Datasets}
\label{sec:datasets}

\subsection{Two distinct data roles}

The datasets serve two different purposes, and the requirements differ:

\begin{enumerate}[leftmargin=2em]
  \item \textbf{Predictor training pool: all tasks, short and long.} The stability label
        is local. It depends only on the recent frames and the local regime, so episode
        length is irrelevant to the label's meaning. Short tasks add cheap regime
        diversity, and they break a shortcut the predictor could otherwise learn, namely
        that late in an episode implies contact.
  \item \textbf{Hypothesis-test tasks: long-horizon only.} The claim under test is that
        no constant $k$ works when a task contains both regimes. Only tasks with both
        regimes can show this. Single-regime short tasks cannot dissociate the
        hypothesis, but they make clean controls: a mostly-stable task should favor long
        chunks or tie, and a contact-heavy task should favor $k=1$ or tie.
\end{enumerate}

\subsection{The specific datasets}

The datasets are organized by the role they play. The rule that governs the split:
the predictor training pool feeds the predictor only, and the held-out evaluation
tasks never enter training.

\paragraph{Predictor training pool (Panda arm with OSC control, one shared predictor).}
\begin{itemize}[leftmargin=2em]
  \item robomimic \textbf{Lift-ph}, \textbf{Can-ph}, \textbf{Square-ph},
        \textbf{Tool-Hang-ph}: all four 200-demo proficient-human sets.
  \item \textbf{MimicGen D0--D1} generated sets across its task suite, including the
        short tasks (coffee, threading, stack) for contact-event diversity.
  \item \textbf{FurnitureSim} \texttt{one\_leg} demonstrations.
\end{itemize}

\paragraph{Label-quality evaluation (calibration and generalization, not policy rollout).}
\begin{itemize}[leftmargin=2em]
  \item Calibration against sim ground-truth FTLE: \textbf{Square-ph} and
        \textbf{Tool-Hang-ph}.
  \item Mostly-stable control task: \textbf{Can-ph}.
  \item Task-transfer test: head trained on Square only, labels scored on Tool-Hang.
        This tests whether the video model learned stability itself or a
        task-specific visual cue.
  \item Embodiment-transfer tier (stage 2): \textbf{ALOHA} simulated insertion and
        transfer-cube, with a separate action head and labels recomputed from
        scratch. Temporal ensembling stays off by default because it reinstates
        per-step feedback; turning it on and off is an ablation.
\end{itemize}

\paragraph{Policy rollout, short horizon.}
\begin{itemize}[leftmargin=2em]
  \item \textbf{Square-ph} and \textbf{Tool-Hang-ph} (primary; reference Diffusion
        Policy numbers).
  \item \textbf{Can-ph} as the negative control: the adaptive schedule should show
        approximately zero gap over the best fixed chunk length here.
\end{itemize}

\paragraph{Policy rollout, long horizon.}
\begin{itemize}[leftmargin=2em]
  \item \textbf{MimicGen} \texttt{three\_piece\_assembly}, \texttt{nut\_assembly},
        \texttt{kitchen}, \texttt{coffee\_preparation}: these span an
        alternation-count axis (how many stable-to-unstable transitions per episode)
        and produce the headline scaling curve.
  \item \textbf{FurnitureSim} \texttt{one\_leg}: the contact-heavy showcase.
  \item \textbf{LIBERO-Long}, paired with \textbf{LIBERO-Pro} perturbation levels: a
        reviewer-legible benchmark with measured headroom.
\end{itemize}

\paragraph{Excluded, deliberately.}
\begin{itemize}[leftmargin=2em]
  \item robomimic \textbf{Transport}: bimanual, demoted to stage 2.
  \item \textbf{CALVIN}: excluded on the TTIC audit consistency findings.
\end{itemize}

\paragraph{Priority order.} Same-embodiment (Panda/OSC) work comes first: the
robomimic ph tasks (running), then MimicGen, then FurnitureSim. LIBERO and ALOHA
follow after the same-embodiment results land.

\paragraph{Real robot (deferred).} Same predictor architecture; labels come from
world-model rollout divergence as described in Section \ref{sec:models}, or from
hardware perturb-and-replay on a subsample.

%=====================================================================
\section{Training}
\label{sec:training}

\subsection{What is trained and what is reused}

\begin{center}
\small
\begin{tabular}{@{}lll@{}}
\toprule
Component & Source & Status \\
\midrule
Chunking policy & Diffusion Policy (real-stanford), ACT & Reused; retrain from \\
 & & released config only if needed \\
Video encoder & V-JEPA 2 ViT-L, frozen & Reused, frozen \\
Stability head & attentive probe + MLP + proprio MLP & \textbf{Built and trained (only one)} \\
Label generator & MuJoCo perturb-and-replay script & Built (a script, not a model) \\
Benchmarks and demos & robomimic, ALOHA sim, FurnitureSim & Reused \\
Real-world label source & V-JEPA 2-AC or DINO-WM divergence & Deferred \\
\bottomrule
\end{tabular}
\end{center}

\subsection{Loss function for the stability head}

\begin{equation}
L \;=\; \mathrm{BCE}_w\bigl(p_{\mathrm{ol}}, y_{\mathrm{ol}}\bigr)
\;+\; \alpha\,\mathrm{Huber}\bigl(\hat\lambda_{\mathrm{ol}}, \lambda_{\mathrm{ol}}\bigr)
\;+\; \beta\,\bigl(\text{the same pair of terms for CL}\bigr).
\end{equation}

\begin{itemize}[leftmargin=2em]
  \item \textbf{Class-weighted BCE} on the open-loop regime is the decision loss.
        Unstable timesteps are the minority class, so we weight the loss or balance the
        sampling.
  \item \textbf{Huber on $\lambda$} is auxiliary. It calibrates the head, and it
        provides the continuous-$k$ rule ($k \approx c/|\hat\lambda|$) for free as a
        later ablation. We start with $\alpha = 0.5$.
  \item \textbf{Closed-loop terms} are trained from the start with $\beta = 0.5$. They
        enable the full three-case runtime rule and the mechanism analysis, at
        negligible cost.
  \item \textbf{Policy losses are untouched.} Diffusion Policy keeps its DDPM MSE loss
        and ACT keeps its L1 plus KL loss. The encoder stays frozen throughout.
\end{itemize}

\subsection{Training procedure and cost}

The encoder is frozen, so encoder features for every clip can be computed once and
cached. Training the head on cached features takes minutes to hours. Policy training,
when needed at all, is at most about 12 hours per task on a single A100 using the
released configurations, and released checkpoints avoid even that. Label generation is
CPU hours. Stage 0 runs on a laptop. Stages 0 and 1 together are estimated at 4 to 6
weeks wall-clock.

%=====================================================================
\section{Experiments}
\label{sec:experiments}

The experiment ladder produces one decisive plot per stage.

\subsection{Stage 0: synthetic sanity check (days, no GPU)}

A synthetic piecewise-linear system with three segments of known contraction rates,
taken from the caveat section of the hypothesis note. We sweep chunk length per segment.
Prediction: the optimal chunk-length profile tracks the contraction profile, with a
floor $\kstar_i$ in contracting segments and a ceiling near $1/\lambda_j$ in expansive
segments, and success degrades as $k \to 1$ in stable segments. This validates the
segment-wise application of the theory before any robot experiment.

\subsection{Stage 1a: the key plot (oracle labels, no predictor)}

A forced-$k$-per-segment sweep on Square and Tool-Hang using the oracle labels directly.
The predictor is not in the loop yet, which isolates the hypothesis from predictor
quality. We plot a heatmap of task success over the grid
$(k_{\mathrm{stable}}, k_{\mathrm{unstable}})$.

Prediction: the best corner is (large, small), and the $k_{\mathrm{stable}} = 1$ column
is \emph{bad}. The second half of this prediction is the sign prediction that no other
method in the literature makes, because every policy-side signal treats $k=1$ as
universally safe.

\subsection{Stage 1b: predictor quality}

AUROC and calibration of $p(\text{OL-unstable})$ against held-out oracle labels;
alignment with simulator contact flags; cross-task generalization (train on Square,
evaluate on Tool-Hang); and the video-versus-single-frame ablation (V-JEPA 2 against
DINOv2 with the same head).

\subsection{Stage 1c: the full closed loop}

The predictor-driven switch against every fixed $k$ and against the policy-side
baselines of Section \ref{sec:baselines}, all running the same policy weights. This
stage includes the two dissociation cases where environment-side and policy-side
signals disagree:

\begin{enumerate}[leftmargin=2em]
  \item \textbf{Contact with clean unimodal visuals.} The policy is confident, but the
        environment is unstable. Entropy and variance baselines over-commit; our
        predictor shortens the chunk.
  \item \textbf{Free space with induced bimodal demonstrations.} The policy is
        uncertain, but the environment is stable. The baselines shorten the chunk, and
        by the theory this per-step replanning is what fails; our predictor commits to
        a long chunk.
\end{enumerate}

\subsection{Stage 2: replication and ablations}

ACT on ALOHA; a long-horizon task (Transport or FurnitureSim); and the ablations:
the continuous-$k$ rule, temporal ensembling on and off, noise injection into
unstable-segment demonstrations per Theorem 2, and the label hyperparameters
$N$, $K$, $\sigma_u$, $\delta$.

%=====================================================================
\section{Baselines}
\label{sec:baselines}

All baselines are inference-time rules on the same trained Diffusion Policy, so every
comparison holds the policy fixed:

\begin{itemize}[leftmargin=2em]
  \item \textbf{Fixed $k \in \{1, 2, 4, 8, 16\}$.} The full constant-chunk sweep.
  \item \textbf{AAC-style.} Execute the low-entropy prefix of each chunk. The Diffusion
        Policy proxy for entropy is the variance across resampled chunks.
  \item \textbf{DVAC-style.} Truncate the chunk when denoising variance rises.
  \item \textbf{BID-style.} Per-step resampling.
\end{itemize}

%=====================================================================
\section{Risks and mitigations}

\begin{center}
\begin{tabular}{@{}p{0.44\linewidth}p{0.5\linewidth}@{}}
\toprule
Risk & Mitigation \\
\midrule
Label noise near segment boundaries &
Deadband $\delta$, median filter, subsample stride. \\
\addlinespace
Perturbation scale wrong: too small hits the noise floor, too large leaves the linear
regime &
$\sigma_u$ set to the policy validation RMSE, plus a sensitivity ablation. \\
\addlinespace
Class imbalance: unstable is the minority class &
Weighted BCE or balanced sampling. \\
\addlinespace
Predictor learns a task-position shortcut (late in episode implies contact) instead of
dynamics &
Cross-task generalization test; short tasks in the training pool; shuffle-clip control. \\
\addlinespace
ACT temporal ensembling reinstates per-step feedback and confounds the comparison &
Off by default, ablated explicitly. \\
\addlinespace
Simulator labels may not transfer to the real world &
World-model-divergence labels replace the oracle later; same head, same losses. \\
\bottomrule
\end{tabular}
\end{center}

%=====================================================================
\section{Compute}

A single A100 or less is sufficient. Diffusion Policy training is at most about 12
hours per task, or zero when released checkpoints exist. Label generation is CPU hours
and embarrassingly parallel. Head training is minutes to hours on cached frozen
features. Stage 0 is a laptop experiment. Stages 0 and 1 are estimated at 4 to 6 weeks
wall-clock.

%=====================================================================
\section{Code and checkpoints}

\begin{itemize}[leftmargin=2em]
  \item Diffusion Policy: \url{github.com/real-stanford/diffusion_policy}
        (robomimic image configs).
  \item robomimic datasets: \url{robomimic.github.io}
        (square, tool\_hang, can; ph variants).
  \item V-JEPA 2, attentive-probe code, and the AC variant for later:\\
        \url{github.com/facebookresearch/vjepa2}.
  \item DINOv2 for the single-frame ablation: \texttt{facebook/dinov2-base}.
  \item ACT and ALOHA sim: \url{github.com/tonyzhaozh/act}.
  \item $\pi_0$ for the real-robot stage:
        \url{github.com/Physical-Intelligence/openpi}.
\end{itemize}

%=====================================================================
\section{Status and results as of July 20, 2026}
\label{sec:status}

This plan was written before execution. Everything through Stage 1 has now run,
and Stage 2 is running. This section records, stage by stage, what was
executed, what was found, and where execution deviated from the plan above. It
is a summary layer; the full evidence lives in four companion documents:
the labeling report (\texttt{label\_report\_latex}), the Stage 1a report
(\texttt{stage1a\_report\_latex}), the predictor document
(\texttt{predictor\_doc\_latex}), and the Stage 2 document
(\texttt{stage2\_doc\_latex}).

\subsection{Stage 0: validated, including the falsifiability direction}

The synthetic piecewise-linear system behaved exactly as Section
\ref{sec:experiments} predicted. Contracting segments showed a floor near
$\kstar \approx 5$ to $6$ (chunks of 3 or less drove success to zero), expansive
segments showed a ceiling near 4 to 5 (chunks of 8 or more drove success to
zero), and with those parameters the best constant $k$ reached 3 percent
success while the adaptive profile reached 100 percent. The FTLE labeler
recovered the true $\lambda$ to four decimals, and the closed-loop $\lambda$
flipped sign correctly in both regimes. A second run with milder parameters
showed a constant $k$ sufficing, which documents the falsifiable boundary:
adaptive chunking wins exactly when the stable-segment floor exceeds the
unstable-segment ceiling, and not otherwise.

\subsection{Labeling: complete on all planned channels}

The open-loop perturb-and-replay procedure of Section \ref{sec:labels} ran to
completion: about 155{,}000 labeled stamps across the four robomimic tasks
(demonstrations plus policy rollouts) and LIBERO. The closed-loop channel ran
on all four robomimic tasks (3{,}969 stamps) and confirmed the mechanism
premise on every task: perturbations injected under closed-loop replanning keep
growing, with positive closed-loop $\lambda$ at 90 to 99 percent of stamps per
task. One planning error is worth recording: the plan called labeling ``CPU
hours''; in practice it was GPU days, with the closed-loop tool\_hang channel
alone taking 54 hours.

\subsection{Stage 1a: the effect exists, and it has a direction}

The forced-$k$ grid ran on square and tool\_hang, 12 cells, three seeds, 150
episodes per cell, with simulator contact plus hysteresis as the oracle
signal. Square is flat (0.773 to 0.873, nothing clears two sigma): the base
policy evaluates at 0.90 and there is no headroom, which is itself the lesson
that Stage 2 task selection acted on. Tool\_hang produced the project's first
rollout-level effect: per-step replanning during contact pools to 0.60 to 0.67
against a 0.79 to 0.81 plateau for committed chunks, at three to five sigma in
three independent cells. The plan's sign prediction, that the $k{=}1$ column
would be bad, held (fixed $k{=}1$ reaches 0.667 on tool\_hang against 0.807
for the best fixed $k$). The surprise is where the damage concentrates: in
contact, the segment the oracle calls unstable. That is the direction the
closed-loop labeling channel predicted (closed-loop $\lambda$ stays positive
in contact, so per-step feedback does not stabilize these segments), and it
contradicts the naive reading that unstable segments simply want $k{=}1$.
The binary oracle switch itself never beat the
best fixed $k$; the explanation is that these rollouts are 65 to 88 percent
contact, so a binary signal leaves the switch almost nothing to modulate. That
finding strengthened, rather than weakened, the case for the continuous
predicted $\lambda$.

\subsection{Stage 1b: the head works, but not for the reason we assumed}

The pooled head reaches AUROC 0.87 on confident held-out stamps (0.89 with
LIBERO in the pool) with $\lambda$ MAE 0.027. Two planned tests came back with
answers the plan did not anticipate. Cross-task generalization fails outright
(AUROC 0.56 training on one task and evaluating on another), so pooled
training is required, and the pooled result is what everything downstream
uses. And the planned video-versus-single-frame ablation returned a tie: a
single-frame DINOv2 head matches the 16-frame V-JEPA 2 head almost exactly
(0.703/0.859 against 0.699/0.868), so the signal the head reads is scene
configuration plus proprioception, not motion. The shuffle-clip control from
the risks table became moot and was dropped: a single frame cannot exploit
temporal order, yet it ties the clip. A head trained on the closed-loop labels
(beyond the plan) reached only 0.627, consistent with closed-loop expansion
being a property of the policy's reaction rather than of the visible scene.

\subsection{Stage 1c: a tie on tasks where a tie is the ceiling}

The predictor-driven $(16, 4)$ switch pooled over three seeds: tool\_hang
0.800, a statistical tie with both the best fixed $k$ (0.807) and the best
oracle cell (0.807), reached while replanning adaptively (mean executed $k$
near 11) and avoiding the per-step trap. Square pools to 0.767 inside the flat
band, the expected null at ceiling. A live perturb-and-replay probe cell was
run once as an in-the-loop oracle but degenerated to fixed $k{=}16$ because
its threshold scale does not transfer from the labeling protocol; it is
recorded and de-prioritized. The Stage 1 conclusion, stated plainly: on short
contact-dominated tasks the best fixed $k$ is near-optimal and matching it is
the achievable outcome; the decisive test of adaptivity needs long tasks that
alternate between stable transport and unstable contact. That conclusion
redefined Stage 2.

\subsection{Stage 2: redefined and now running}

The Stage 2 of Section \ref{sec:experiments} (ACT on ALOHA, execution-rule
ablations) is deferred, not dropped. In its place, Stage 2 is the long-horizon
test: Diffusion Policy on four MimicGen tasks (three\_piece\_assembly,
nut\_assembly, kitchen, coffee\_preparation, D0, 1{,}000 demonstrations each,
mean episode lengths 337 to 689 steps), followed by the same fixed-$k$ versus
predictor-switch comparison. The first training runs were killed at 80 to 86
percent of their 2,000-epoch budget by a scratch-disk failure early on
July 21; checkpoints saved every 50 epochs to persistent storage, so
selection proceeds on the surviving grids (which end past the quality
plateau) while fresh full-length reruns provide insurance in parallel. The
experiment timeline moved by less than a day: checkpoint selection the
evening of July 21, the switch-versus-fixed-$k$ rollouts on July 22. The
design, the simulator version conflict that shapes the infrastructure, the
disk incident, and the zero-shot decision for the head are documented in the
Stage 2 companion document.

\subsection{Deviations and still-open items}

\begin{itemize}[leftmargin=2em]
  \item \textbf{Compute.} The plan's estimates were wrong in both directions
        of kind: Diffusion Policy training is about five to six days per task
        at the 200k-step budget, not 12 hours, and labeling was GPU days, not
        CPU hours. The work ran on an 8-GPU machine, not a single A100.
  \item \textbf{DINOv2 variant.} The ablation used \texttt{dinov2-large}, not
        the \texttt{dinov2-base} listed in the code and checkpoints section.
  \item \textbf{Not yet run.} The policy-side baselines (AAC, DVAC, BID
        styles), the two dissociation cases of Stage 1c, the can negative
        control, the continuous-$k$ rule, the temporal-ensembling and
        noise-injection ablations, and the label-hyperparameter sweeps. These
        remain valid follow-ups; none blocks Stage 2.
  \item \textbf{Parked.} FurnitureSim (manual download requirement);
        LIBERO-Long policy training is an open option for currently idle
        GPUs.
\end{itemize}

%=====================================================================
\section{Summary}

The project trains exactly one small network: a head on a frozen video encoder that
predicts whether the current moment is open-loop stable or unstable. The labels come
from a simulator procedure that nudges the first action by the size of error the policy
actually makes and fits the slope of the log divergence. The labels themselves define
the segments; nothing is hand-segmented. At run time the predictor switches the executed
chunk length: long chunks in stable segments, per-step feedback in unstable ones. The
policy is never modified, so every comparison isolates the execution rule. The central
testable prediction, which no policy-confidence method can make, is that short chunks
actively fail in stable segments. Stage 1a tested exactly that with oracle labels before
the predictor was ever in the loop; Section \ref{sec:status} records the outcome,
including the three-to-five-sigma failure of per-step replanning during contact on
tool\_hang.

\end{document}
